
\chapter{Memory and Storage}
\label{chap:memory-and-storage}

\chapterintro{This chapter explains the difference between temporary memory and long-term storage, and why memory limits matter when working with LLMs.}
\learningobjectives{\begin{itemize}\item Distinguish RAM from storage.\item Explain cache, memory hierarchy, and latency.\item Understand why LLMs require large amounts of memory.\end{itemize}}

\section{RAM}
Random Access Memory stores data that programs are actively using. RAM is fast, but temporary.

\section{Storage}
SSDs and hard drives store data for the long term. They are slower than RAM but retain information after power is removed.

\section{Memory Hierarchy}
\begin{center}\begin{tabular}{ll}\toprule Level & Description\\\midrule Registers & Tiny and extremely fast\\ Cache & Very fast CPU memory\\ RAM & Main working memory\\ SSD & Long-term storage\\ Network Storage & Remote storage\\\bottomrule\end{tabular}\end{center}

\section{LLM Memory Requirements}
Large models contain billions of parameters. Each parameter is a number. Storing and using these numbers requires memory. Quantization reduces memory by using fewer bits per parameter.

\section{Python Example}
\begin{lstlisting}
import sys
numbers = [1, 2, 3, 4, 5]
print(sys.getsizeof(numbers))
\end{lstlisting}

\section{Exercise}
\begin{exercise}Why might a model fail to load on a computer with limited RAM or GPU memory?\end{exercise}
\begin{solution}The model parameters, intermediate activations, and runtime overhead may require more memory than the device has available.\end{solution}

\section{Summary}
Memory and storage determine what data and models a system can handle. Production AI engineering requires careful attention to memory, especially for embeddings, vector stores, and LLM inference.
